\chapter{Couleur globale}
\section{Température et nuance}
La température couvre 2\,000 à 50\,000 K ; la nuance couvre $-150$ à $+150$.
La température déplace l’équilibre chaud/froid ; la nuance corrige l’axe
vert/magenta. Les bandes du curseur sont des indications, pas un nuancier
contractuel du résultat. La balance automatique se trouve dans le groupe RAW.

La température de référence dépend du développement de la photo ; elle n’est
pas toujours 5\,200 K ou 6\,500 K. La remise à zéro du contrôle de température
revient à sa température de référence, pas à zéro kelvin. La nuance se
réinitialise à zéro dans l’interface actuelle. Les unités ou valeurs affichées
par un autre logiciel ne garantissent pas une correspondance numérique exacte.

\section{Vibrance et saturation}
Les deux contrôles couvrent $-100$ à $+100$ avec une référence à zéro.
La saturation modifie la richesse colorée globale. La vibrance cherche une
action plus sélective suivant les couleurs déjà présentes ; ce n’est pas
simplement une saturation moins forte. Les effets peuvent devenir importants
sur les couleurs vives et modifier leur luminosité apparente.

\paragraph{Exemple.}
Commencer par une petite correction de température, puis régler la nuance.
Ajouter ensuite une vibrance modérée pour les couleurs ternes et vérifier
les carnations ainsi que les rouges et jaunes. N’ajuster la saturation
globale qu’en regardant l’ensemble de l’image et les zones déjà colorées.
Maintenir \touche{Espace} dans la vue pour comparer au rendu d’origine.

\attention{Les validations de couleur dépendent du rendu de départ, du profil
et du boîtier. Les valeurs extrêmes sont utiles pour tester, pas une promesse
de reproduction exacte de Lightroom.}
\gestesreglages
